\subsection{Umbral Recurrences and Variants}

\begin{theorem} \label{thm:associated-sequence-recurrence-and-transfer} \hlt{Recurrence and Transfer Formulas}\\
Let $(p_n)$ be associated with the delta series $f$, and set $u(t)\defeq f(t)/t\in\mathscr{F}^{\times}$. Then:
\begin{enumerate}[label=(\roman*)]
\setlength{\itemsep}{0pt}
\item $p_{n+1}=x[f'(t)]^{-1}p_n$ for $n\geq0$.
\item $p_n=f'(t)u(t)^{-n-1}x^n=xu(t)^{-n}x^{n-1}$ for $n\geq1$.
\end{enumerate}
\begin{proof}
\begin{enumerate}[label=(\roman*)]
\setlength{\itemsep}{0pt}
\item By Theorem \ref{thm:umbral-shift-derivation-chain-rule}, $\theta_t=(\partial_t f)\theta_f=f'(t)\theta_f$. Since $\theta_t=x$ and $\theta_fp_n=p_{n+1}$, $p_{n+1}=x[f'(t)]^{-1}p_n$.
\item Put $q_n\defeq f'u^{-n-1}x^n$. Then $q_0=f'u^{-1}1=1=p_0$ and, using $xa(t)=a(t)x-a'(t)$,\\
$q_{n+1}=f'u^{-n-2}x^{n+1}=xu^{-n-1}x^n=x(f')^{-1}q_n$. Thus $q_n=p_n$ by (i).\\
The same commutator gives $xu^{-n}x^{n-1}=u^{-n}x^n+nu^{-n-1}u'x^{n-1}=f'u^{-n-1}x^n$.
\end{enumerate}
\hlt{Structural Commentary}: The recurrence determines $(p_n)$; transfer solves it using $f=tu$.\\
\hlt{Key Insight}: The Pincherle commutator converts the shift recurrence into closed form.
\end{proof}
\end{theorem}

\begin{corollary} \label{cor:sheffer-shift-formula} \hlt{Sheffer Shift Formula}\\
Let $(s_n)$ be Sheffer for $(g,f)$, with $f$ delta and $g\in\mathscr{F}^{\times}$. Write $f'=f'(t)$ and $g'=g'(t)$. Then:
\begin{enumerate}[label=(\roman*)]
\setlength{\itemsep}{0pt}
\item $\theta_{g,f}=g^{-1}\theta_fg=\left(x-\frac{g'}{g}\right)(f')^{-1}$.
\item $s_{n+1}=\left(x-\frac{g'}{g}\right)(f')^{-1}s_n$ for $n\geq0$.
\end{enumerate}
\begin{proof}
\begin{enumerate}[label=(\roman*)]
\setlength{\itemsep}{0pt}
\item Since $gs_n$ is associated with $f$, $\theta_{g,f}=g^{-1}\theta_fg$. By Theorem \ref{thm:associated-sequence-recurrence-and-transfer}, $\theta_f=x(f')^{-1}$.\\
Since $f'$ and $g$ commute, and $g'=gx-xg$, $g^{-1}\theta_fg=g^{-1}x(f')^{-1}g=g^{-1}xg(f')^{-1}=\left(x-\frac{g'}{g}\right)(f')^{-1}$.
\item By definition of the Sheffer shift, $\theta_{g,f}s_n=s_{n+1}$ for $n\geq0$. Using (i), $s_{n+1}=\theta_{g,f}s_n=\left(x-\frac{g'}{g}\right)(f')^{-1}s_n$.
\end{enumerate}
\hlt{Structural Commentary}: Sheffer shifts are conjugated associated shifts.\\
\hlt{Key Insight}: Conjugation by $g$ contributes the correction $g'/g$.
\end{proof}
\end{corollary}

\begin{method} \label{meth:building-sheffer-sequences} \hlt{Building Sheffer Sequences}\\
To construct the Sheffer sequence for $(g,f)$:
\begin{enumerate}[label=(\roman*)]
\setlength{\itemsep}{0pt}
\item Verify $o(f)=1$ and $o(g)=0$.
\item If $\overline f$ is explicit, compute $S(y,t)\defeq g(\overline f(t))^{-1}e^{y\overline f(t)}$, where $s_n(y)=n![t^n]S(y,t)$.
\item If $\overline f$ is unavailable, solve recursively from $f(t)s_n=ns_{n-1}$, where $\langle g(t)f(t)^k\mid s_n\rangle=n!\delta_{n,k}$.
\item For associated sequences, set $g=1$ and use Theorem
\ref{thm:associated-sequence-recurrence-and-transfer}.
\end{enumerate}
\end{method}

\begin{definition} \label{def:stirling-and-exponential-polynomials} \hlt{Stirling and Exponential Polynomials}
\begin{enumerate}[label=(\roman*)]
\setlength{\itemsep}{0pt}
\item \hlt{Stirling numbers}: $S(n,k)$ are defined by $x^n=\sum\limits_{k=0}^{n}S(n,k)(x)_k$.
\item \hlt{Exponential polynomials}: define $\phi_n(x)\defeq\sum\limits_{k=0}^{n}S(n,k)x^k$.
\item \hlt{Boundary convention}: $S(n,-1)=0$ and $S(n,k)=0$ for $k>n$.
\end{enumerate}
\end{definition}

\begin{example} \label{ex:umbral-recurrences-stirling} \hlt{Lower Factorials and Stirling Numbers}\\
For $f(t)=e^t-1$, $f'(t)=e^t$, hence $p_{n+1}(x)=xe^{-t}p_n(x)=xp_n(x-1)$. With $p_0=1$, this gives $p_n=(x)_n$.\\
For $f(t)=\log(1+t)$, $\overline f(t)=e^t-1$, and the associated sequence is $(\phi_n)$.\\
The inverse umbral relation is $x^n=\sum\limits_{k=0}^{n}S(n,k)(x)_k$, and $\phi_n(x)=\sum\limits_{k=0}^{n}S(n,k)x^k$.\\
The recurrence formula gives $\phi_{n+1}(x)=x(1+t)\phi_n(x)=x\bigl(\phi_n(x)+\phi_n'(x)\bigr)$,\\
so coefficient comparison yields $S(n+1,k)=S(n,k-1)+kS(n,k)$.
\end{example}

\begin{definition} \label{def:q-and-logarithmic-umbral-variants} \hlt{\texorpdfstring{$q$}{q}- and Logarithmic Umbral Variants}
\begin{enumerate}[label=(\roman*)]
\setlength{\itemsep}{0pt}
\item \hlt{Gaussian coefficient}: for $0\leq k\leq n$,
$\binom{n}{k}_q\defeq
\prod\limits_{j=1}^{n}(1-q^j)/
\left(\prod\limits_{j=1}^{k}(1-q^j)
\prod\limits_{j=1}^{n-k}(1-q^j)\right)$.
\item \hlt{\texorpdfstring{$q$}{q}-umbral calculus}: the deformation $\binom{n}{k}\leadsto\binom{n}{k}_q$.
\item \hlt{Logarithmic umbral calculus}: the logarithmic-type variant based on harmonic logarithms.
\end{enumerate}
\end{definition}

\begin{remark} \label{rem:umbral-calculus-further-directions} \hlt{Further Directions}\\
Classical umbral methods extend to many polynomial families and operations:
\begin{enumerate}[label=(\roman*)]
\setlength{\itemsep}{0pt}
\item \hlt{Families}: lower factorial, Stirling, exponential, Gould, central
factorial, Abel, Mittag--Leffler, Bessel, Bell, Hermite, Bernoulli, Euler,
Laguerre, Poisson--Charlier, actuarial, Meixner, Pidduck, Narumi, Boole,
Peters, Mahler, and Mott polynomials.
\item \hlt{Operations}: connection constants, duplication formulas, Lagrange
inversion, cross sequences, Steffensen sequences, operational formulas,
inverse relations, recurrence solving, and binomial convolution.
\item \hlt{Variants}: $q$-umbral and logarithmic umbral calculi replace or
modify the binomial structure through the objects in Definition
\ref{def:q-and-logarithmic-umbral-variants}.
\end{enumerate}
\end{remark}
